
\documentclass{article}
\usepackage{colortbl}
\usepackage{makecell}
\usepackage{multirow}
\usepackage{supertabular}

\begin{document}

\newcounter{utterance}

\centering \large Interaction Transcript for game `cladder', experiment `full\_v1.5\_default', episode 6429 with Qwen3.6{-}27B.
\vspace{24pt}

{ \footnotesize  \setcounter{utterance}{1}
\setlength{\tabcolsep}{0pt}
\begin{supertabular}{c@{$\;$}|p{.15\linewidth}@{}p{.15\linewidth}p{.15\linewidth}p{.15\linewidth}p{.15\linewidth}p{.15\linewidth}}
    \# & \multicolumn{2}{c}{Player} && \multicolumn{2}{c}{Game Master} \\
    \hline

    \theutterance \stepcounter{utterance}  
    & & & \multicolumn{4}{p{0.6\linewidth}}{
        \cellcolor[rgb]{0.9,0.9,0.9}{
            \makecell[{{p{\linewidth}}}]{
                \texttt{\tiny{[P1$\langle$GM]}}
                \texttt{You are an expert in causal inference. The following question is not a typical commonsense query, but rather a meticulously designed question created by a professor specializing in causal inference, intended to assess the students' mastery of the course content.} \\
\texttt{Imagine a self{-}contained, hypothetical world with only the following conditions, and without any unmentioned factors or causal relationships: Rixq has a direct effect on xevu. Zuph has a direct effect on xevu. The overall probability of rixq is 40\%. The probability of not rixq and xevu is 9\%. The probability of rixq and xevu is 22\%. Is the chance of xevu smaller when observing rixq?} \\
\texttt{Start your answer with 'Yes' or 'No', followed by additional reasoning or evidence to support your explanation.} \\
            }
        }
    }
     \\ \\

    \theutterance \stepcounter{utterance}  
    & \multicolumn{4}{p{0.6\linewidth}}{
        \cellcolor[rgb]{0.561,0.737,0.561}{
            \makecell[{{p{\linewidth}}}]{
                \texttt{\tiny{[P1$\rangle$GM]}}
                \texttt{No} \\
\\ 
\texttt{To determine whether the chance of `xevu` is smaller when observing `rixq`, we need to compare the conditional probability \$P(xevu | rixq)\$ with the marginal probability \$P(xevu)\$.} \\
\\ 
\texttt{**Step 1: Calculate \$P(xevu | rixq)\$**} \\
\texttt{The formula for conditional probability is:} \\
\texttt{\$\$P(xevu | rixq) = \textbackslash{}frac\{P(rixq \textbackslash{}cap xevu)\}\{P(rixq)\}\$\$} \\
\\ 
\texttt{Given:} \\
\texttt{*   \$P(rixq \textbackslash{}cap xevu) = 22\textbackslash{}\% = 0.22\$} \\
\texttt{*   \$P(rixq) = 40\textbackslash{}\% = 0.40\$} \\
\\ 
\texttt{\$\$P(xevu | rixq) = \textbackslash{}frac\{0.22\}\{0.40\} = 0.55\$\$} \\
\\ 
\texttt{So, the probability of `xevu` given `rixq` is **55\%**.} \\
\\ 
\texttt{**Step 2: Calculate the marginal probability \$P(xevu)\$**} \\
\texttt{The total probability of `xevu` is the sum of the probability of `xevu` occurring with `rixq` and the probability of `xevu` occurring without `rixq`.} \\
\\ 
\texttt{\$\$P(xevu) = P(rixq \textbackslash{}cap xevu) + P(\textbackslash{}neg rixq \textbackslash{}cap} \\
            }
        }
    }
    & & \\ \\

    \theutterance \stepcounter{utterance}  
    & & & \multicolumn{2}{p{0.3\linewidth}}{
        \cellcolor[rgb]{0.9,0.9,0.9}{
            \makecell[{{p{\linewidth}}}]{
                \texttt{\tiny{[GM$|$GM]}}
                \texttt{No} \\
            }
        }
    }
    & & \\ \\

    \theutterance \stepcounter{utterance}  
    & & & \multicolumn{2}{p{0.3\linewidth}}{
        \cellcolor[rgb]{0.9,0.9,0.9}{
            \makecell[{{p{\linewidth}}}]{
                \texttt{\tiny{[GM$|$GM]}}
                \texttt{no} \\
            }
        }
    }
    & & \\ \\

    \theutterance \stepcounter{utterance}  
    & & & \multicolumn{2}{p{0.3\linewidth}}{
        \cellcolor[rgb]{0.9,0.9,0.9}{
            \makecell[{{p{\linewidth}}}]{
                \texttt{\tiny{[GM$|$GM]}}
                \texttt{game\_result = WIN} \\
            }
        }
    }
    & & \\ \\

\end{supertabular}
}

\end{document}
